% ================================================================
% CHAPTER 2: RESONANCE AS INNER PRODUCT
% ================================================================
\chapter{Resonance as Inner Product}
\label{ch:resonance}

\epigraph{A sign is something by knowing which we know something more.}{Charles Sanders Peirce, \textit{Collected Papers}, 8.332}

\section{From Binary Resonance to Quantitative Resonance}

In IDT~v1, resonance was a binary relation: two ideas either resonate ($a \res
b$) or they do not. This captured the \emph{existence} of intellectual affinity
but not its \emph{degree}. The passage to IDT~v2 replaces this binary relation
with a real-valued function --- the \textbf{resonance strength} --- defined as
the inner product of the embedded vectors.

\begin{definition}[Resonance Strength]
\label{def:resonance-strength}
The \textbf{resonance strength} between ideas $a$ and $b$ is:
\[
  \rs{a}{b} \;:=\; \ip{\embed{a}}{\embed{b}}_\Hspace.
\]
\end{definition}

This definition inherits all the properties of the inner product, which we now
develop systematically.


\section{Fundamental Properties of Resonance}

\begin{theorem}[Symmetry of Resonance]
\label{thm:resonance-symmetry}
For all ideas $a, b \in \IS$:
\[
  \rs{a}{b} = \rs{b}{a}.
\]
\end{theorem}

\begin{proof}
This follows immediately from the symmetry of the real inner product:
\[
  \rs{a}{b} = \ip{\embed{a}}{\embed{b}} = \ip{\embed{b}}{\embed{a}} = \rs{b}{a}.
\]
In Lean~4:
\begin{lstlisting}
theorem rs_symm (a b : Idea) : rs a b = rs b a := by
  simp [rs, real_inner_comm]
\end{lstlisting}
\end{proof}

\begin{interpretation}[Mutuality of Resonance]
\label{interp:mutuality}
The symmetry of resonance is a mathematical expression of a deep philosophical
principle: \emph{intellectual affinity is mutual}. If idea $a$ resonates with
idea $b$, then $b$ resonates equally with $a$. Darwin's theory resonates with
Wallace's just as much as Wallace's resonates with Darwin's. This is not
trivially true in all frameworks: in a directed model of ``influence,'' the
relationship might be asymmetric (Darwin influenced Wallace more than Wallace
influenced Darwin). But resonance, as we define it, is not influence --- it
is \emph{structural affinity}, and structural affinity is inherently symmetric.

In Bourdieu's sociology, this corresponds to the observation that cultural
capital is not merely possessed but \emph{recognized}: agent $A$'s recognition
of $B$'s cultural standing is necessarily reciprocated by $B$'s recognition of
$A$'s, at least at the structural level. The asymmetries of power and prestige
are modeled not by asymmetric resonance but by differences in \emph{norm} (some
ideas carry more weight than others) and by differences in \emph{strategic
position} (some agents occupy more favorable positions in the communication
game).
\end{interpretation}

\begin{theorem}[Self-Resonance]
\label{thm:self-resonance}
For any idea $a \in \IS$:
\[
  \rs{a}{a} = \nrm{\embed{a}}^2 \geq 0,
\]
with equality if and only if $\embed{a} = 0$ (i.e., $a$ is synonymous with
$\void$ in the embedding).
\end{theorem}

\begin{proof}
$\rs{a}{a} = \ip{\embed{a}}{\embed{a}} = \nrm{\embed{a}}^2$ by definition of
the norm. Non-negativity and the characterization of equality follow from the
axioms of an inner product space.

In Lean~4:
\begin{lstlisting}
theorem rs_self_nonneg (a : Idea) : 0 ≤ rs a a := by
  simp [rs, inner_self_nonneg]

theorem rs_self_eq_zero_iff (a : Idea) :
    rs a a = 0 ↔ φ a = 0 := by
  simp [rs, inner_self_eq_zero]
\end{lstlisting}
\end{proof}

\begin{definition}[Norm of an Idea]
\label{def:idea-norm}
The \textbf{norm} (or \textbf{weight}) of an idea $a$ is:
\[
  \nrm{a}_\IS \;:=\; \nrm{\embed{a}}_\Hspace = \sqrt{\rs{a}{a}}.
\]
\end{definition}

The norm measures the ``magnitude'' or ``weight'' of an idea in meaning space.
It is always non-negative, and it equals zero only for ideas that are synonymous
with the void.

\begin{theorem}[Void Has Zero Resonance]
\label{thm:void-zero-resonance}
For all ideas $a \in \IS$:
\[
  \rs{\void}{a} = 0 \quad \text{and} \quad \rs{a}{\void} = 0.
\]
\end{theorem}

\begin{proof}
$\rs{\void}{a} = \ip{\embed{\void}}{\embed{a}} = \ip{0}{\embed{a}} = 0$.
The second equality follows from symmetry, or directly by the same argument
applied to the second component.

In Lean~4:
\begin{lstlisting}
theorem rs_void (a : Idea) : rs void a = 0 := by
  simp [rs, embed_void, inner_zero_left]
\end{lstlisting}
\end{proof}

\begin{interpretation}[Silence Resonates with Nothing]
\label{interp:silence-resonance}
The void has zero resonance with every idea. This is the mathematical expression
of a profound truth about silence: genuine silence --- not the pregnant pause or
the strategic withholding of speech, but the utter absence of ideation --- has
no intellectual affinity with anything. It neither agrees nor disagrees,
neither supports nor opposes, neither reinforces nor undermines.

This should be contrasted with \emph{near-void} ideas: ideas of very small norm,
which are close to silence but not identical with it. A vague platitude, a
meaningless pleasantry, a formulaic greeting --- these have small norm and
therefore small resonance with most other ideas, but they are not exactly the
void. They occupy the neighborhood of the origin in meaning space: the region
of near-silence, where ideas barely exist as ideas at all.

In the Buddhist tradition, this corresponds to the distinction between
\emph{\'s\=unyat\=a} (emptiness, the void) and ordinary ignorance. The void is
not mere absence of knowledge but a positive state of non-ideation --- the zero
vector, not a small random vector.
\end{interpretation}


\section{The Cauchy--Schwarz Inequality: Limits of Understanding}

The most important inequality in all of functional analysis is the
Cauchy--Schwarz inequality. In the context of IDT~v2, it becomes a fundamental
limit on the degree of resonance between ideas.

\begin{theorem}[Cauchy--Schwarz for Resonance]
\label{thm:cauchy-schwarz}
For all ideas $a, b \in \IS$:
\[
  |\rs{a}{b}|^2 \;\leq\; \rs{a}{a} \cdot \rs{b}{b},
\]
or equivalently:
\[
  |\rs{a}{b}| \;\leq\; \nrm{\embed{a}} \cdot \nrm{\embed{b}}.
\]
Equality holds if and only if $\embed{a}$ and $\embed{b}$ are linearly dependent
--- i.e., one is a scalar multiple of the other.
\end{theorem}

\begin{proof}
This is the Cauchy--Schwarz inequality for the inner product space $\Hspace$,
applied to the vectors $\embed{a}$ and $\embed{b}$. The proof is standard:
for any $t \in \mathbb{R}$,
\[
  0 \leq \nrm{\embed{a} + t\embed{b}}^2
  = \nrm{\embed{a}}^2 + 2t\ip{\embed{a}}{\embed{b}} + t^2\nrm{\embed{b}}^2.
\]
This is a non-negative quadratic in $t$, so its discriminant is non-positive:
\[
  4\ip{\embed{a}}{\embed{b}}^2 - 4\nrm{\embed{a}}^2\nrm{\embed{b}}^2 \leq 0,
\]
which gives the result. Equality holds iff the quadratic has a real root, i.e.,
iff $\embed{a} + t\embed{b} = 0$ for some $t$, i.e., iff $\embed{a}$ and
$\embed{b}$ are linearly dependent.

In Lean~4:
\begin{lstlisting}
theorem cauchy_schwarz_resonance (a b : Idea) :
    rs a b ^ 2 ≤ rs a a * rs b b := by
  simp [rs]
  exact inner_mul_le_norm_mul_sq (φ a) (φ b)
\end{lstlisting}
\end{proof}

\begin{interpretation}[The Fundamental Limit on Cross-Cultural Understanding]
\label{interp:cauchy-schwarz}
The Cauchy--Schwarz inequality for resonance is arguably the single most
important result in IDT~v2. It states that the degree of resonance between
two ideas is bounded above by the product of their individual weights. In
normalized form:
\[
  \frac{|\rs{a}{b}|}{\sqrt{\rs{a}{a}} \cdot \sqrt{\rs{b}{b}}} \;\leq\; 1.
\]
This ratio --- the absolute value of the \emph{cosine similarity} between
$\embed{a}$ and $\embed{b}$ --- measures the ``purity'' of the resonance:
how much of each idea's weight is involved in the resonance with the other.
The bound says that this purity can never exceed $1$, and it reaches $1$
only when the ideas are ``parallel'' --- pointing in the same (or exactly
opposite) direction in meaning space.

In cultural terms, this is a mathematical limit on cross-cultural understanding.
No matter how empathetic, how scholarly, how committed to understanding another
culture you are, the degree to which your ideas can resonate with theirs is
bounded by the product of your respective cultural weights. And this bound is
tight only when your ideas are ``parallel'' to theirs --- when you are, in some
sense, already thinking in the same direction.

This formalizes a version of Kuhn's thesis of paradigm incommensurability, but
with an important refinement. Kuhn claimed that scientists working in different
paradigms cannot fully understand each other. The Cauchy--Schwarz inequality
agrees --- but it quantifies the degree of mutual understanding, rather than
treating incommensurability as an all-or-nothing phenomenon. Two paradigms can
have small but non-zero resonance, corresponding to partial but incomplete
mutual understanding. And the bound tells us exactly how much resonance is
possible given the ``weights'' of the two paradigms.

The equality case is also illuminating. Maximum resonance occurs when
$\embed{a} = \lambda \embed{b}$ for some scalar $\lambda$. In cultural terms,
this means that full mutual understanding requires the ideas to be ``parallel''
--- they must point in the same direction in meaning space, differing only in
magnitude. This is the formal counterpart of the hermeneutic insight that
understanding requires a ``fusion of horizons'' (Gadamer): you can fully
understand another's idea only when your own idea is already oriented in the
same direction.
\end{interpretation}


\section{Normalized Resonance and Cosine Similarity}

The Cauchy--Schwarz inequality motivates the following normalization.

\begin{definition}[Normalized Resonance]
\label{def:normalized-resonance}
For ideas $a, b \in \IS$ with $\embed{a} \neq 0$ and $\embed{b} \neq 0$,
the \textbf{normalized resonance} (or \textbf{cosine similarity}) is:
\[
  \hat{r}(a, b) \;:=\; \frac{\rs{a}{b}}{\nrm{\embed{a}} \cdot \nrm{\embed{b}}}
  \;=\; \frac{\ip{\embed{a}}{\embed{b}}}{\nrm{\embed{a}} \cdot \nrm{\embed{b}}}.
\]
\end{definition}

\begin{proposition}[Range of Normalized Resonance]
\label{prop:normalized-range}
For non-void ideas $a$ and $b$:
\[
  -1 \;\leq\; \hat{r}(a, b) \;\leq\; 1.
\]
Moreover:
\begin{enumerate}
\item $\hat{r}(a, b) = 1$ if and only if $\embed{a}$ and $\embed{b}$ are
positive scalar multiples of each other (``perfect alignment'');
\item $\hat{r}(a, b) = -1$ if and only if $\embed{a}$ and $\embed{b}$ are
negative scalar multiples of each other (``perfect opposition'');
\item $\hat{r}(a, b) = 0$ if and only if $\embed{a} \perp \embed{b}$
(``complete independence'').
\end{enumerate}
\end{proposition}

\begin{proof}
The bounds follow from the Cauchy--Schwarz inequality (Theorem~\ref{thm:cauchy-schwarz}).
The characterizations of the extreme cases follow from the equality conditions
of Cauchy--Schwarz and the definition of orthogonality.
\end{proof}

\begin{interpretation}[The Resonance Spectrum]
\label{interp:spectrum}
The normalized resonance $\hat{r}(a, b) \in [-1, 1]$ provides a complete
spectrum of intellectual relationships:

\begin{center}
\begin{tabular}{cl}
\textbf{Value} & \textbf{Interpretation} \\
\hline
$+1$ & Perfect alignment: ideas say the ``same thing'' (up to emphasis) \\
$(0, 1)$ & Partial agreement: ideas point in roughly the same direction \\
$0$ & Independence: ideas are about completely unrelated topics \\
$(-1, 0)$ & Partial conflict: ideas point in roughly opposite directions \\
$-1$ & Perfect opposition: ideas are exactly contradictory \\
\end{tabular}
\end{center}

This spectrum is richer than both the binary resonance of IDT~v1 (which
conflates all positive values into ``resonates'' and treats zero and negative
values as ``does not resonate'') and the informal language of intellectual
discourse (which typically distinguishes only between ``agreement,''
``disagreement,'' and ``irrelevance'').

The cosine similarity formulation also connects IDT~v2 to the vast literature
on distributional semantics and information retrieval, where cosine similarity
is the standard measure of semantic relatedness between text vectors. In that
literature, the cosine similarity between two document vectors measures how much
they are ``about the same thing.'' IDT~v2 elevates this computational technique
to a principled theoretical framework.
\end{interpretation}


\section{Bilinearity and the Algebra of Resonance}

The inner product is bilinear: linear in each argument separately. This gives
resonance a rich algebraic structure.

\begin{theorem}[Bilinearity of Resonance]
\label{thm:bilinearity}
For all ideas $a, b, c \in \IS$:
\begin{enumerate}
\item $\rs{a \compose b}{c} = \rs{a}{c} + \rs{b}{c}$ \hfill
(left-additivity)
\item $\rs{c}{a \compose b} = \rs{c}{a} + \rs{c}{b}$ \hfill
(right-additivity)
\end{enumerate}
\end{theorem}

\begin{proof}
For (1):
\begin{align*}
  \rs{a \compose b}{c}
  &= \ip{\embed{a \compose b}}{\embed{c}} \\
  &= \ip{\embed{a} + \embed{b}}{\embed{c}} \\
  &= \ip{\embed{a}}{\embed{c}} + \ip{\embed{b}}{\embed{c}} \\
  &= \rs{a}{c} + \rs{b}{c}.
\end{align*}
Part (2) follows by symmetry, or by the same argument using linearity in the
second component.

In Lean~4:
\begin{lstlisting}
theorem rs_bilinear_left (a b c : Idea) :
    rs (a ∘ b) c = rs a c + rs b c := by
  simp [rs, embed_compose, inner_add_left]
\end{lstlisting}
\end{proof}

\begin{corollary}[Resonance of Compositions]
\label{cor:resonance-compositions}
For ideas $a, b, c, d \in \IS$:
\[
  \rs{a \compose b}{c \compose d}
  = \rs{a}{c} + \rs{a}{d} + \rs{b}{c} + \rs{b}{d}.
\]
\end{corollary}

\begin{proof}
Apply left-additivity, then right-additivity to each term:
\begin{align*}
  \rs{a \compose b}{c \compose d}
  &= \rs{a}{c \compose d} + \rs{b}{c \compose d} \\
  &= \rs{a}{c} + \rs{a}{d} + \rs{b}{c} + \rs{b}{d}.
\end{align*}
\end{proof}

\begin{interpretation}[The Resonance Decomposition]
\label{interp:decomposition}
Corollary~\ref{cor:resonance-compositions} says that the resonance between two
composite ideas decomposes into a sum of ``pairwise'' resonances between their
components. The resonance between ``democratic socialism'' and ``market
liberalism'' is the sum of four terms: democracy--markets, democracy--liberalism,
socialism--markets, and socialism--liberalism. Each pair contributes independently
to the total resonance.

This decomposition has important implications for \emph{coalition politics}.
When two political ideologies are composites of several component ideas, their
mutual resonance depends on all pairwise interactions between components. A
coalition is stable when the sum is positive (overall resonance) and unstable
when it is negative (overall conflict). The decomposition shows exactly which
pairwise interactions are contributing to or detracting from the coalition.

More abstractly, the bilinearity of resonance means that the resonance function
$\rs{\cdot}{\cdot}$ is a \emph{bilinear form} on the image of $\varphi$. This
places IDT~v2 squarely within the classical theory of quadratic forms, with all
the algebraic machinery that entails: diagonalization, signature, index, and
Witt equivalence. The spectral theory developed in Chapter~\ref{ch:spectral}
exploits this connection systematically.
\end{interpretation}


\section{Resonance and the Riesz Representation Theorem}

The Riesz representation theorem states that every continuous linear functional
on a Hilbert space is represented by inner product with a unique vector. In
IDT~v2, this has a beautiful interpretation.

\begin{proposition}[Ideas as Linear Functionals]
\label{prop:riesz}
Each idea $c \in \IS$ defines a linear functional $\ell_c : \Hspace \to
\mathbb{R}$ by $\ell_c(v) := \ip{v}{\embed{c}}$. By the Riesz representation
theorem, $\ell_c$ is the unique continuous linear functional represented by the
vector $\embed{c}$.
\end{proposition}

\begin{proof}
The map $v \mapsto \ip{v}{\embed{c}}$ is clearly linear and continuous (by
Cauchy--Schwarz). The Riesz representation theorem asserts that this is the
\emph{only} way to define a continuous linear functional on $\Hspace$: every
such functional has the form $v \mapsto \ip{v}{w}$ for a unique $w \in \Hspace$.
\end{proof}

\begin{interpretation}[Ideas as Evaluation Functions]
\label{interp:riesz}
The Riesz representation establishes a profound duality: an idea is
simultaneously a \emph{point} in meaning space (the vector $\embed{c}$) and an
\emph{evaluation function} on meaning space (the functional $\ell_c$). As a
point, the idea has a location, a direction, and a magnitude. As a functional,
the idea evaluates other meanings, assigning each a real-valued ``score'' ---
the resonance strength.

This duality is deeply Peircean. In Peirce's semiotics, a sign has three
aspects: the \emph{representamen} (the sign-vehicle), the \emph{object} (what
the sign refers to), and the \emph{interpretant} (the effect of the sign on an
interpreter). The vector $\embed{c}$ corresponds to the representamen (the sign
as a mathematical object); the functional $\ell_c$ corresponds to the
interpretant (the sign's capacity to evaluate and be evaluated by other signs).
The object is the ``external reality'' to which the sign refers, which IDT~v2
does not model directly --- it models only the relationships \emph{among} signs.

The Riesz representation also connects to Bourdieu's theory of cultural capital.
Bourdieu argued that cultural objects function simultaneously as ``things'' (with
intrinsic properties) and as ``instruments of distinction'' (used to evaluate
and classify other agents). An artwork is both an object in a gallery and a
criterion of taste. A philosophical theory is both a set of propositions and a
standard against which other theories are judged. The point-functional duality
of the Hilbert space model captures this double character precisely.
\end{interpretation}


\section{Peirce's Semiotics and the Three Types of Resonance}

Charles Sanders Peirce distinguished three types of signs based on the
relationship between the sign and its object:

\begin{itemize}
\item \textbf{Icons}: signs that resemble their objects (a portrait, a map, a
diagram);
\item \textbf{Indices}: signs that are causally connected to their objects
(smoke as a sign of fire, a weathervane as a sign of wind direction);
\item \textbf{Symbols}: signs whose relationship to their objects is purely
conventional (words, mathematical symbols, traffic lights).
\end{itemize}

In IDT~v2, these three types of signs correspond to three sources of resonance.

\begin{definition}[Iconic, Indexical, and Symbolic Resonance]
\label{def:peirce-resonance}
Let $a, b \in \IS$ be ideas with $\rs{a}{b} \neq 0$. We say the resonance is:
\begin{itemize}
\item \textbf{Iconic} if $\embed{a}$ and $\embed{b}$ are close in direction
(small angle), reflecting structural similarity in meaning;
\item \textbf{Indexical} if the resonance arises from ideas that co-occur
in the same contexts (occupying similar regions of $\Hspace$ due to shared
contextual features);
\item \textbf{Symbolic} if the resonance is the result of conventional
association --- the embedding $\varphi$ maps them to correlated vectors by
cultural convention rather than structural necessity.
\end{itemize}
\end{definition}

\begin{interpretation}[The Sources of Meaning]
\label{interp:peirce}
This classification reveals that the inner product $\rs{a}{b}$ aggregates
multiple sources of meaning-relatedness into a single number. Iconic resonance
reflects inherent structural similarity (``democracy'' resonates with
``republic'' because the concepts share structural features). Indexical resonance
reflects contextual co-occurrence (``smoke'' resonates with ``fire'' because
they appear in the same situations). Symbolic resonance reflects arbitrary
convention (``dog'' resonates with ``loyalty'' in English-language cultural
context, but this resonance is conventional, not structural or causal).

The Hilbert space model does not distinguish among these sources --- they all
contribute to the same inner product. This is both a strength (it provides a
unified measure of meaning-relatedness) and a limitation (it cannot tell us
\emph{why} two ideas resonate). Future extensions of IDT~v2 might decompose
$\Hspace$ into ``iconic,'' ``indexical,'' and ``symbolic'' subspaces, with
the total resonance being the sum of projections onto each.

The Peircean trichotomy also suggests that the embedding $\varphi$ is not
unique: different embedding choices emphasize different sources of resonance.
A ``structural'' embedding (emphasizing logical relationships) would maximize
iconic resonance; a ``distributional'' embedding (emphasizing co-occurrence)
would maximize indexical resonance; a ``cultural'' embedding (emphasizing
conventional associations) would maximize symbolic resonance. Each embedding
captures a different aspect of meaning, and the full picture requires
considering multiple embeddings simultaneously --- a theme we return to in
Chapter~\ref{ch:spectral}.
\end{interpretation}


\section{The Polarization Identity}

The inner product and the norm are related by the \emph{polarization identity},
which allows us to express resonance in terms of distances.

\begin{theorem}[Polarization Identity for Resonance]
\label{thm:polarization}
For all ideas $a, b \in \IS$:
\[
  \rs{a}{b} = \frac{1}{2}\left(
    \nrm{\embed{a} + \embed{b}}^2 - \nrm{\embed{a}}^2 - \nrm{\embed{b}}^2
  \right)
  = \frac{1}{2}\left(
    \rs{a \compose b}{a \compose b} - \rs{a}{a} - \rs{b}{b}
  \right).
\]
\end{theorem}

\begin{proof}
Expanding:
\begin{align*}
  \nrm{\embed{a} + \embed{b}}^2
  &= \ip{\embed{a} + \embed{b}}{\embed{a} + \embed{b}} \\
  &= \nrm{\embed{a}}^2 + 2\ip{\embed{a}}{\embed{b}} + \nrm{\embed{b}}^2 \\
  &= \nrm{\embed{a}}^2 + 2\rs{a}{b} + \nrm{\embed{b}}^2.
\end{align*}
Solving for $\rs{a}{b}$ gives the result. The second form uses $\nrm{\embed{a}
+ \embed{b}}^2 = \nrm{\embed{a \compose b}}^2 = \rs{a \compose b}{a \compose
b}$.

In Lean~4:
\begin{lstlisting}
theorem polarization_resonance (a b : Idea) :
    rs a b = (rs (a ∘ b) (a ∘ b) - rs a a - rs b b) / 2 := by
  simp [rs, embed_compose, inner_add_left, inner_add_right]
  ring
\end{lstlisting}
\end{proof}

\begin{interpretation}[Resonance from Observables]
\label{interp:polarization}
The polarization identity has a remarkable practical consequence: it shows that
resonance between two ideas can be \emph{computed} from three self-resonance
measurements. If you can measure the ``weight'' of idea $a$ alone, the weight
of idea $b$ alone, and the weight of their composition $a \compose b$, then you
can compute their resonance:
\[
  \rs{a}{b} = \frac{\rs{a \compose b}{a \compose b} - \rs{a}{a} - \rs{b}{b}}{2}.
\]
This is analogous to measuring the gravitational force between two masses by
observing their combined motion rather than directly detecting the force.

In empirical applications, self-resonance might be measured as the ``cultural
weight'' or ``salience'' of an idea --- how much attention it commands, how
frequently it is discussed, how strongly people feel about it. The polarization
identity then says: the resonance between two ideas can be inferred from the
salience of each idea separately and the salience of their combination. If
combining two ideas produces a composite with more salience than the sum of the
parts, the ideas have positive resonance (they reinforce each other). If the
composite has less salience, the resonance is negative (they partially cancel).

This provides a potentially \emph{testable} prediction of IDT~v2: the
``cultural salience'' of a composite idea should differ from the sum of the
component saliences by exactly twice the resonance. Any systematic deviation
from this prediction would indicate a failure of the additivity axiom --- a
non-linearity in how ideas combine.
\end{interpretation}


\section{Self-Resonance as Conviction}

Self-resonance deserves special attention, as it plays a central role in the
game theory of meaning (Part~II).

\begin{definition}[Conviction]
\label{def:conviction}
The \textbf{conviction} of an idea $a$ is its self-resonance:
\[
  \kappa(a) \;:=\; \rs{a}{a} = \nrm{\embed{a}}^2.
\]
\end{definition}

\begin{proposition}[Properties of Conviction]
\label{prop:conviction}
Conviction satisfies:
\begin{enumerate}
\item $\kappa(a) \geq 0$ for all $a$, with $\kappa(a) = 0$ iff $a
\equiv_\varphi \void$.
\item $\kappa(a \compose b) = \kappa(a) + \kappa(b) + 2\rs{a}{b}$.
\item $\kappa(a^{\langle n \rangle}) = n^2 \kappa(a)$.
\end{enumerate}
\end{proposition}

\begin{proof}
(1) follows from the positivity of the inner product. For (2):
\begin{align*}
  \kappa(a \compose b)
  &= \nrm{\embed{a} + \embed{b}}^2 \\
  &= \nrm{\embed{a}}^2 + 2\ip{\embed{a}}{\embed{b}} + \nrm{\embed{b}}^2 \\
  &= \kappa(a) + 2\rs{a}{b} + \kappa(b).
\end{align*}
(3) follows from Corollary~\ref{cor:self-resonance-iterated} in
Chapter~\ref{ch:geometry}.
\end{proof}

\begin{interpretation}[Conviction and Composition]
\label{interp:conviction}
The formula $\kappa(a \compose b) = \kappa(a) + \kappa(b) + 2\rs{a}{b}$ reveals
a non-trivial interaction between conviction and resonance. When you combine two
ideas that have positive resonance ($\rs{a}{b} > 0$), the conviction of the
composite exceeds the sum of the individual convictions. Reinforcing ideas
produce more conviction than they carry separately --- a phenomenon familiar
from political rhetoric, where combining multiple aligned messages produces a
``bandwagon effect'' that exceeds what any individual message would achieve.

Conversely, when the resonance is negative ($\rs{a}{b} < 0$), the composite has
\emph{less} conviction than the sum of the parts. This is the formal expression
of \emph{cognitive dissonance}: holding two conflicting beliefs simultaneously
reduces one's overall conviction. If the conflict is severe enough
($\rs{a}{b} < -(\kappa(a) + \kappa(b))/2$, which is impossible by
Cauchy--Schwarz unless the norms are equal and the ideas are anti-parallel),
the composite conviction approaches zero.

The quadratic growth $\kappa(a^{\langle n \rangle}) = n^2 \kappa(a)$ means that
conviction grows super-linearly with repetition. This is the ``echo chamber''
effect: an idea repeated within an insular community gains disproportionate
conviction, not because of any evidence or argument, but simply because the
vector gets longer with each repetition, and conviction is the \emph{square}
of the length.
\end{interpretation}

\begin{example}[Resonance in Political Discourse]
\label{ex:political-resonance}
Consider three political ideas embedded in $\mathbb{R}^3$:
\begin{align*}
  \embed{\text{equality}} &= (3, 0, 1), \\
  \embed{\text{liberty}} &= (2, 2, 0), \\
  \embed{\text{order}} &= (0, -1, 3).
\end{align*}
The resonance matrix is:
\[
  \begin{pmatrix}
    \rs{\text{eq}}{\text{eq}} & \rs{\text{eq}}{\text{lib}} &
    \rs{\text{eq}}{\text{ord}} \\
    \rs{\text{lib}}{\text{eq}} & \rs{\text{lib}}{\text{lib}} &
    \rs{\text{lib}}{\text{ord}} \\
    \rs{\text{ord}}{\text{eq}} & \rs{\text{ord}}{\text{lib}} &
    \rs{\text{ord}}{\text{ord}}
  \end{pmatrix}
  =
  \begin{pmatrix}
    10 & 6 & 3 \\
    6 & 8 & -2 \\
    3 & -2 & 10
  \end{pmatrix}.
\]
Equality and liberty resonate positively (6), but liberty and order have
negative resonance ($-2$): they are in mild tension. Equality and order have
moderate positive resonance (3). The ideology ``liberty + order'' has conviction
$8 + 10 + 2(-2) = 14$, while ``equality + liberty'' has conviction
$10 + 8 + 2(6) = 30$. The reinforcing pair produces more than twice the
conviction of the conflicting pair.
\end{example}

\begin{example}[Cosine Similarity in Distributional Semantics]
\label{ex:cosine}
In computational linguistics, word2vec embeddings yield cosine similarities such
as:
\begin{align*}
  \hat{r}(\text{king}, \text{queen}) &\approx 0.73, \\
  \hat{r}(\text{king}, \text{woman}) &\approx 0.32, \\
  \hat{r}(\text{king}, \text{potato}) &\approx 0.04.
\end{align*}
The IDT~v2 framework provides a principled interpretation of these numbers:
king and queen are highly resonant (closely aligned in meaning space), king and
woman are moderately resonant (sharing some semantic features but differing in
others), and king and potato are nearly orthogonal (semantically unrelated). The
inner product gives a quantitative measure of what distributional semantics
discovers empirically.
\end{example}

\bigskip

This chapter has established resonance as the inner product of embedded ideas,
developed its fundamental properties, and connected it to semiotics, cultural
theory, and distributional semantics. In Chapter~\ref{ch:metric}, we derive the
metric structure that the inner product induces on the space of ideas --- the
\emph{distance} between ideas, which will give precise mathematical content to
notions of intellectual proximity and paradigmatic incommensurability.
